% SGA 3, Expose VI_B, Loop 1 English reconstruction.
% Scope: Corollary 2.8 through Corollary 2.12 in full;
% Section 3 and paragraph 3.0 are excluded.
% Authority: Polo--Gille Exp6B-13oct24.pdf, component-local pages 14--16
% (combined-reader pages 355--357; original-margin printed pages 337--340).
% OCR and the Reinhold English Markdown are comparison witnesses only.

\medskip
\phantomsection\label{cor:sga3-VIB-2.8}%
\noindent\textbf{推论 2.8.---}%
\footnote{\label{note:sga3-VIB-21}编者注: 将``拟有限''情形的断言
与``非分歧''情形的断言分开, 改变了表述方式.}
\emph{设 \(u:G\to H\) 为 \(S\)-群态射.
假设对每个 \(s\in S\), \(G_s\) 在 \(\kappa(s)\) 上局部有限型.}
\begin{enumerate}
\item[\textup{a)}] \emph{若 \(u\) 局部有限型, 则以下条件等价:}
  \begin{enumerate}
  \item[\textup{(i)}] \emph{\(u\) 局部拟有限;}
  \item[\textup{\SGAInlineRefTarget{cor:sga3-VIB-2.8:item:ii}(ii)}]
    \emph{对每个 \(s\in S\), 态射 \(u_s:G_s\to H_s\) 局部拟有限;}
  \item[\textup{\SGAInlineRefTarget{cor:sga3-VIB-2.8:item:iii}(iii)}]
    \emph{\(\operatorname{Ker}u\) 在 \(S\) 上局部拟有限;}
  \item[\textup{(iv)}] \emph{\(\operatorname{Ker}u\) 的纤维离散.}
  \end{enumerate}
\item[\textup{b)}] \emph{若 \(u\) 局部有限呈示, 则以下条件等价:}
  \begin{enumerate}
  \item[\textup{(v)}] \emph{\(u\) 非分歧;}
  \item[\textup{(vi)}] \emph{对每个 \(s\in S\), 态射 \(u_s:G_s\to H_s\) 非分歧;}
  \item[\textup{\SGAInlineRefTarget{cor:sga3-VIB-2.8:item:vii}(vii)}]
    \emph{\(\operatorname{Ker}u\) 非分歧;}
  \item[\textup{(viii)}] \emph{单位截面 \(S\to\operatorname{Ker}u\) 为开浸入.}
  \end{enumerate}
\end{enumerate}

等价关系
\[
  \textup{(i)}\Longleftrightarrow\textup{(ii)}
  \qquad\text{以及}\qquad
  \textup{(v)}\Longleftrightarrow\textup{(vi)}
\]
来自 \hyperref[concl:sga3-VIB-2.5-ii]{2.5}\textup{\SGARefLink{cor:sga3-VIB-2.8:item:ii}{(ii)}},
下面的回顾~\hyperref[rem:sga3-VIB-2.8.1]{2.8.1} 表明
\[
  \textup{(i)}\Longrightarrow\textup{(iii)}
  \qquad\text{以及}\qquad
  \textup{(v)}\Longrightarrow\textup{(vii)}.
\]

对每个 \(s\in S\), 用 \(e_s\) 表示 \(H_s\) 的单位元.
则 \textup{\SGARefLink{cor:sga3-VIB-2.8:item:iii}{(iii)}}
(分别地, \textup{\SGARefLink{cor:sga3-VIB-2.8:item:vii}{(vii)}})
蕴含对每个 \(s\in S\),
\[
  (\operatorname{Ker}u)_s
    =\operatorname{Ker}u_s
    =u_s^{-1}(e_s)
\]
在 \(\kappa(s)=\kappa(e_s)\) 上局部拟有限 (分别地, 非分歧).
由此可知 \(u_s\) 在 \(G_s\) 的单位点处拟有限 (分别地, 非分歧),
因而由 \SGARefLink{prop:sga3-VIB-1.3}{1.3},
\(u_s\) 局部拟有限 (分别地, 非分歧). 故
\[
  \textup{(iii)}\Longrightarrow\textup{(ii)}
  \qquad\text{以及}\qquad
  \textup{(vii)}\Longrightarrow\textup{(vi)}.
\]
最后,
\[
  \textup{(ii)}\Longleftrightarrow\textup{(iv)}
  \quad\text{由 1.4.1,}\qquad
  \textup{(vii)}\Longleftrightarrow\textup{(viii)}
  \quad\text{由 2.7.}
\]

\medskip
\phantomsection\label{rem:sga3-VIB-2.8.1}%
\noindent\textbf{回顾 \SGARefLink{rem:sga3-VIB-2.8.1}{2.8.1}.---}
回顾一下 (参见 \SGARefLink{sga3:i:definition:2-3-6-1}{第~I~篇, 2.3.6.1}),
给定 \(S\)-群态射 \(u:G\to H\), \emph{\(u\) 的核},
记为 \(\operatorname{Ker}u\), 是 \(G\) 的群子函子,
对每个态射 \(f:T\to S\) 定义为
\[
  (\operatorname{Ker}u)(T)
    =\{a\in G(T)\mid u\circ a=\epsilon_H\circ f\}.
\]
由上述文献 (或 EGA~I, 4.4.1), 该函子由 \(S\)-群
\[
  G\times_H S=u^{-1}\bigl(\epsilon_H(S)\bigr),
\]
表示, 后者也简称为 \(\operatorname{Ker}u\).
特别地, 结构态射
\[
  \operatorname{Ker}u\longrightarrow S
\]
由 \(u\) 作基变换得到.

\medskip
\phantomsection\label{lem:sga3-VIB-2.9}%
\noindent\textbf{引理 2.9.---}
\emph{设 \(\pi:X\to S\) 是具有 \(S\)-截面 \(\epsilon\) 的态射.}
\begin{enumerate}
\item[\textup{(i)}] \emph{若 \(\pi\) 为单射, 则它是整态射.}%
  \begingroup
  \renewcommand{\thefootnote}{\ensuremath{(\ast)}}%
  \footnote[1]{\label{note:sga3-VIB-2.9-star}这也是 EGA~IV\(_4\), 18.12.11
  的一个特例, 因为 \(\pi\) 显然为普遍同胚.}%
  \endgroup
\item[\textup{(ii)}] \emph{若 \(\pi\) 局部有限型,
  且对每个 \(s\in S\), \(\pi_s\) 为同构, 则 \(\pi\) 为同构.}%
  \begingroup
  \renewcommand{\thefootnote}{\ensuremath{(\ast\ast)}}%
  \footnote[2]{\label{note:sga3-VIB-2.9-starstar}这也是 EGA~IV\(_4\), 18.12.6
  的直接推论.}%
  \endgroup
\end{enumerate}

先注意到, 由下面的 \hyperref[lem:sga3-VIB-2.9.1]{引理~2.9.1},
由于 \(\pi\) 是同胚, 它是仿射态射.

若 \(\pi\) 为单射, 则 \(\epsilon\) 为满射.
由于 \(\epsilon\) 是满射浸入, \(\epsilon(S)\) 同构于
\(X\) 中由 \(\mathcal O_X\) 的一个由幂零元组成的理想 \(\mathcal I\) 定义的闭子概形.
由于 \(\epsilon\) 是 \(\pi\) 的 \(S\)-截面,
有 \(\mathcal O_S\)-模的直和分解
\[
  \mathcal O_X=\mathcal O_S\oplus\mathcal I
\]
由于 \(\mathcal I\) 是 \(\mathcal O_X\) 中由幂零元组成的理想,
它显然在 \(\mathcal O_S\) 上整.
因而 \(\mathcal O_X\) 在 \(\mathcal O_S\) 上整, 故 \(\pi\) 为整态射.

现假设 \(\pi\) 局部有限型.
由于 \(\pi\circ\epsilon=\operatorname{id}_S\),
截面 \(\epsilon\) 局部有限呈示 (参见 EGA~IV\(_1\), 1.4.3(v));
因而 \(\mathcal I\) 是 \(\mathcal O_X\) 中的有限型理想
(EGA~IV\(_1\), 1.4.1). 对每个 \(s\in S\), 有
\[
  \mathcal O_{X_s}
    =\kappa(s)\oplus
      \mathcal I\otimes_{\mathcal O_S}\kappa(s).
\]
由假设, \(\epsilon_s\) 为同构, 所以对每个 \(s\in S\),
\[
  \mathcal I\otimes_{\mathcal O_S}\kappa(s)=0
\]
因而对每个 \(x\in X\), 更有
\[
  \mathcal I\otimes_{\mathcal O_X}\kappa(x)=0
\]
现在由 Nakayama 引理, \(\mathcal I=0\), 故 \(\pi\) 为同构.

\medskip
\phantomsection\label{lem:sga3-VIB-2.9.1}%
\noindent\textbf{引理 2.9.1.---}
\emph{设 \(f:X\to Y\) 是为同胚的概形态射. 则 \(f\) 是仿射态射.}%
\begingroup
\renewcommand{\thefootnote}{\ensuremath{(\ast\ast\ast)}}%
\footnote[3]{\label{note:sga3-VIB-2.9.1-starstarstar}
参见 EGA~IV\(_4\), 18.12.7.1, 其中以相同论证证明了一个稍一般的结果.}%
\endgroup

只须说明每个点 \(y\in Y\) 都有开邻域 \(W\),
使 \(f\) 在 \(W\) 上的限制仿射.
设 \(y\in Y\), \(V\) 为 \(y\) 在 \(Y\) 中的仿射开邻域. 令
\[
  V'=f^{-1}(V).
\]
则 \(V'\) 是 \(x=f^{-1}(y)\) 在 \(X\) 中的开邻域.
在 \(X\) 中存在包含于 \(V'\) 的 \(x\) 的仿射开邻域 \(W'\).
令
\[
  W=f(W').
\]
则 \(W\) 是 \(y\) 在 \(Y\) 中的开邻域,
包含在仿射概形 \(V\) 中, 所以 \(W\) 分离.
由于 \(W'\) 仿射, \(f\) 在 \(W\) 上的限制为仿射态射
(EGA~II, 1.2.3).

\medskip
\phantomsection\label{cor:sga3-VIB-2.10}%
\noindent\textbf{推论 2.10.---}
\emph{设 \(G\) 为局部有限型的 \(S\)-群.
假设对每个 \(s\in S\), \(G_s\) 为单位 \(\kappa(s)\)-群.
则 \(G\) 为单位 \(S\)-群.}

更一般地:

\medskip
\phantomsection\label{cor:sga3-VIB-2.11}%
\noindent\textbf{推论 2.11.---}
\emph{设 \(f:X\to Y\) 为局部有限型的 \(S\)-态射.
\(f\) 为单态射的充分必要条件是对每个 \(s\in S\),
\(f_s\) 都为单态射.}

必要性显然; 证明其充分性.
若对每个 \(s\in S\), \(f_s\) 为单态射,
则对每个 \(y\in Y\), 更有 \(f_y\) 为单态射.
因而可假设 \(Y=S\).

由 EGA~I, 5.3.8, 为证明 \(f\) 为单态射, 只须说明
\[
  \Delta_f:X\longrightarrow X\times_SX
\]
是同构, 或等价地, 第一个投影
\[
  p:X\times_SX\longrightarrow X
\]
是同构. 若 \(f_s\) 为单态射, EGA~I, 5.3.8 同样表明,
可等同于 \(p_s\) 的第一个投影
\[
  X_s\times_{\kappa(s)}X_s\longrightarrow X_s,
\]
是同构. 但 \(p\) 具有 \(S\)-截面 \(\Delta_f\);
故 \hyperref[lem:sga3-VIB-2.9]{引理~2.9} 表明,
若对每个 \(s\in S\), \(p_s\) 为同构, 则 \(p\) 本身为同构.

\medskip
\phantomsection\label{cor:sga3-VIB-2.12}%
\noindent\textbf{推论 2.12.---}
\emph{设 \(G\) 为具有 \(S\)-截面 \(\epsilon\) 的 \(S\)-概形,
并假设结构态射 \(\pi:G\to S\) 为闭映射.
设 \(s\in S\), 使 \(\pi\) 在 \(\epsilon(s)\) 处为有限呈示, 且}
\[
  \epsilon_s:\operatorname{Spec}\bigl(\kappa(s)\bigr)
    \longrightarrow G_s
\]
\emph{为同构 (或等价地, 使}
\[
  \pi_s:G_s\longrightarrow
    \operatorname{Spec}\kappa(s)
\]
\emph{为同构). 则 \(s\) 在 \(S\) 中存在开邻域 \(U\), 使}
\[
  \epsilon|_U:U\longrightarrow G\times_SU
\]
\emph{为同构.}

令 \(V\) 为 \(G\) 中使 \(\pi\) 非分歧的点组成的集合.
已知此集合开 (参见编者\hyperref[note:sga3-VIB-6]{注~6}),
且包含 \(\epsilon(s)\). 因而
\[
  W=\epsilon^{-1}(V)
\]
是 \(S\) 中包含 \(s\) 的开子集,
且对每个 \(t\in W\), \(\pi\) 在 \(\epsilon(t)\) 处非分歧.
由于 \(\pi\) 闭, \(\pi|_W\) 也闭, 因而可假设 \(W=S\).

由 \hyperref[lem:sga3-VIB-2.7.1]{2.7.1}, 此时
\[
  X=G-\epsilon(S)
\]
是 \(G\) 的闭子集, 且不与 \(G_s\) 相交.
由于 \(\pi\) 闭,
\[
  F=\pi(X)
\]
是 \(S\) 中不包含 \(s\) 的闭子集. 令
\[
  U=S-F.
\]
则 \(U\) 是 \(S\) 的开子集,
而 \(\epsilon|_U\) 是从 \(U\) 到 \(G|_U\) 的同构.
